\begin{lemma}[Factorial-root product with exceptional rows]
For real numbers $D>0$ and $\theta>0$ with $D+3\theta<1$, there is an
integer $y_0>0$ such that for every integer $y\ge y_0$ and positive
integers $r_0,\ldots,r_{y-1}\le y$, if
\[
 k:=|\{i:r_i>(D+\theta)y\}|\le\theta y,
\]
then
\[
 \prod_{i=0}^{y-1}(r_i!)^{1/r_i}\le(D+3\theta)^y y!.
\]
The threshold depends only on $D$ and $\theta$.
\end{lemma}
\begin{proof}
Put $a_r=(r!)^{1/r}$ for positive integers $r$.
This sequence is increasing: $a_r\le a_{r+1}$ is equivalent, after
raising to the positive power $r(r+1)$ and cancelling $(r!)^r$, to
$r!\le(r+1)^r$, which follows by bounding each factor by $r+1$.
For $1\le r\le y$, comparison of the increasing function $\log t$ with
its integrals over $[1,r]$ gives
\[
 r\log r-r\le\log(r!)\le r\log r-r+\log r+1.
\]
Indeed the lower bound follows from
$\sum_{j=2}^r\log j\ge\int_1^r\log t\,dt=r\log r-r+1$,
and the upper bound follows from
$\sum_{j=1}^{r-1}\log j\le\int_1^r\log t\,dt$ after adding $\log r$.
Consequently
\[
 \frac{a_r}{a_y}\le\frac r y
       \exp\left(\frac{\log y+1}{r}\right).                 \tag{1}
\]

Set $x=D+2\theta$ and
$\gamma=\log((x+\theta)/x)-\theta\log(1/x)$.
We have $0<x<x+\theta<1$ and $\theta<x/2$.
To see $\gamma>0$, integrating $1/t$ over $[x,x+\theta]$ gives
$\log(1+\theta/x)>\theta/(x+\theta)$.
The function $t\log(1/t)$ on $(0,1)$ has maximum $1/e$ (differentiate:
its derivative is $\log(1/t)-1$). Thus
$(x+\theta)\log(1/x)<(3/2)x\log(1/x)\le3/(2e)<1$,
and hence $\theta/(x+\theta)>\theta\log(1/x)$.
The elementary strict inequality $e>3/2$ follows, for example, by
integrating $\exp t\ge1$ over $[0,1]$ to obtain $e\ge2$.

Choose $y_0$ large enough that for all integers $y\ge y_0$,
$y>0$, $1/y\le\theta$, and
\[
 \delta_y:=\frac{\log y+1}{\theta y}\le\gamma/2.             \tag{2}
\]
This is possible since $\log y/y\to0$ and $1/y\to0$.
For completeness the former limit follows from
$\log y=\int_1^y t^{-1}dt\le\int_1^y t^{-1/2}dt
=2(\sqrt y-1)$ for $y\ge1$. These choices use only $D,\theta$.

For an exceptional index, monotonicity gives $a_{r_i}/a_y\le1$.
For a nonexceptional index with $r_i\ge\theta y$, (1) gives
$a_{r_i}/a_y\le(D+\theta)e^{\delta_y}\le xe^{\delta_y}$.
For a remaining index put $s=\lceil\theta y\rceil$.
Then $1\le s$, $r_i\le s$, and
$\theta y\le s\le\theta y+1\le2\theta y<y$,
the last inequality following from $3\theta<1$.
Apply monotonicity and (1) to $s$; since $s/y\le2\theta\le x$,
the same bound $a_{r_i}/a_y\le xe^{\delta_y}$ follows.
Thus all nonexceptional rows satisfy this bound and
\[
 \frac{\prod_i a_{r_i}}{y!}
 \le(xe^{\delta_y})^{y-k}
 =\exp((y-k)\log x+(y-k)\delta_y)
 \le\exp\bigl(y((1-\theta)\log x+\delta_y)\bigr).           \tag{3}
\]
Here $a_y^y=y!$, $\log x<0$, $k\le\theta y$, and
$0\le k\le y$, $\delta_y\ge0$ justify the last inequality even
when $xe^{\delta_y}$ is greater than one.
Finally
\[
 (1-\theta)\log x+\delta_y
 \le\log x+\theta\log(1/x)+\gamma
 =\log(x+\theta).
\]
Substituting this in (3), exponentiating, and multiplying by the
positive number $y!$ proves the claim.
\end{proof}
